% !TeX root = ../main.tex
\chapter{绪论}

\section{研究背景及意义}

地球物理反演通过地表或井中观测到的地震波、重磁、电磁等地球物理场响应，重建地下介质的物性参数空间分布，是深部资源勘探、灾害预警和城市地下空间探测的核心科学问题\cite{tarantola2005inverse, zhao1995mechanics}。随着勘探目标逐渐向深层、隐蔽油气藏以及复杂构造区转移，地下介质呈现出显著的非均匀性、各向异性与多尺度嵌套特征，对反演成像的分辨率与可靠性提出了更高要求\footnote{本模板脚注按规范第 1.12 节采用当页脚注方式，每页自动从带圈数字 \ding{172} 开始重新编号，采用小五号宋体排版。}。

传统反演方法以最小二乘优化与伴随状态法为代表，在理论体系上较为完备，但在面对强非线性、观测孔径受限及含噪数据时存在严重的不适定性\cite{gbt7714, wang2023seismic}。另一方面，伴随深度学习在图像重建与科学计算中的兴起，融合领域先验知识与数据驱动模型的智能反演方法成为当前地球物理与人工智能交叉领域的研究热点\cite{raissi2019physics, zhang2024inversion}。

\section{国内外研究现状}

\subsection{传统确定性地球物理反演方法}

经典线性化反演与全波形反演（Full-Waveform Inversion, FWI）通过最小化合成记录与实测数据之间的残差泛函迭代更新地下模型参数\cite{tarantola2005inverse}。为克服反演多解性，研究者相继引入了 Tikhonov 正则化、全变差（Total Variation）约束以及多尺度频带延拓策略\cite{desmarais1992carbon, yang2021fwi}。尽管上述方法在低频信息充足时能够获得高精度成像结果，但其庞大的正演模拟开销与严苛的初始模型要求限制了实时成像与大规模三维应用。

\subsection{基于深度学习的智能反演方法}

近年来，卷积神经网络、U-Net 及其变体被广泛应用于地震波阻抗反演与速度建模任务\cite{wu2019inversionnet, li2022dl}。然而，纯数据驱动方法在训练样本分布与目标工区存在差异时容易出现预测漂移。为提升物理一致性，物理信息神经网络（Physics-Informed Neural Networks, PINN）将偏微分方程残差直接引入损失函数\cite{raissi2019physics}，为高泛化智能反演提供了新思路。

\section{本文主要研究内容与结构安排}

围绕上述关键科学问题，本文共分为四章，具体组织结构如下：
\begin{enumerate}[label=（\arabic*）]
  \item 第一章为绪论，阐述研究背景、国内外研究现状及本文的总体框架；
  \item 第二章为研究方法与理论分析，建立多尺度特征融合网络与物理约束联合反演模型，并演示公式、图、三线表与跨页续表的规范排法；
  \item 第三章为实验结果与讨论，结合合成数据与实际工区资料开展定量评价与消融分析；
  \item 第四章为结论与展望，总结全文主要创新点，指出当前研究局限并展望未来工作方向。
\end{enumerate}
